\chapter{Contexto e partes interessadas}\label{cap:descricao}

\section{Perspectiva e responsabilidades}
O produto recebe eventos publicados e associa inscrições a uma identidade
externa. O catálogo de origem é responsável por publicar dados válidos; o
provedor de identidade é responsável pela autenticação. O produto aplica as
regras de prazo, capacidade e titularidade ao executar cada operação.

\begin{quadrosrs}
\caption{Necessidades e responsabilidades das partes interessadas}
\begin{tabelasrs}{L{0.23\textwidth}YY}
\cabecalhosrs\textbf{Parte interessada} & \textbf{Necessidade} & \textbf{Responsabilidade no exemplo}\\\midrule
STK-001: visitante & Encontrar eventos futuros. & Consultar o catálogo público.\\
STK-002: participante & Inscrever-se e acompanhar seus registros. & Confirmar ou cancelar a própria inscrição.\\
STK-003: responsável pelo evento & Preservar capacidade e prazo. & Publicar eventos no serviço externo e definir suas condições.\\
STK-004: equipe do produto & Entender e verificar as obrigações. & Manter contratos e relacionar requisitos a evidências de teste.\\
\end{tabelasrs}
\fonte{Elaboração didática deste guia.}
\end{quadrosrs}

\section{Fronteira do sistema}
A figura mostra somente atores e serviços externos. A persistência das
inscrições faz parte da responsabilidade do produto e não é um ator externo.

\begin{figurasrs}
\caption{Contexto do sistema de inscrições}\label{fig:contexto}
\centering
\begin{tikzpicture}[
  caixa/.style={draw=srslinha,rounded corners=2pt,align=center,font=\small,text width=30mm,minimum height=13mm},
  principal/.style={caixa,fill=srscabecalho,text width=36mm,minimum height=23mm},
  seta/.style={-{Stealth[length=2mm]},semithick},
  rotulo/.style={font=\footnotesize,align=center,text width=22mm,inner sep=1pt}
]
\node[principal] (srs) at (0,0) {Sistema de\\inscrições};
\node[caixa] (visitante) at (-6.2,1.6) {Visitante};
\node[caixa] (participante) at (-6.2,-1.6) {Participante};
\node[caixa] (catalogo) at (6.2,1.6) {Catálogo\\externo};
\node[caixa] (identidade) at (6.2,-1.6) {Provedor de\\identidade};
\draw[seta] (visitante.east) -- node[rotulo,above=3mm]{consulta\\pública} (srs.north west);
\draw[seta] (participante.east) -- node[rotulo,below=3mm]{inscrição e\\cancelamento} (srs.south west);
\draw[seta] (catalogo.west) -- node[rotulo,above=3mm]{eventos\\publicados} (srs.north east);
\draw[seta] (identidade.west) -- node[rotulo,below=3mm]{identidade e\\sessão} (srs.south east);
\end{tikzpicture}
\fonte{Elaboração didática deste guia.}
\end{figurasrs}

\Needspace{22\baselineskip}
\section{Ambiente, premissas e dependências}
O ambiente de referência é uma aplicação web em português brasileiro, acessada
por navegador com HTML, CSS e JavaScript. A interface apresenta datas no fuso
\texttt{America/Sao\_Paulo}, enquanto os instantes de referência são mantidos
em UTC. A especificação não impõe framework ou fornecedor de banco de dados.

\begin{quadrosrs}
\caption{Condições que delimitam a operação}
\begin{tabelasrs}{L{0.23\textwidth}YY}
\cabecalhosrs\textbf{Tipo e ID} & \textbf{Condição} & \textbf{Consequência}\\\midrule
Premissa PRE-001 & Relógio do servidor é a referência de prazo. & Comparar início e operação no mesmo referencial.\\
Dependência DEP-001 & Catálogo externo fornece eventos válidos. & Entrada inválida é rejeitada; preservar a última versão válida.\\
Dependência DEP-002 & Provedor externo fornece identidade e sessão. & Novas sessões dependem do provedor; sessões existentes valem até expirar.\\
Restrição RES-001 & Verificação de vaga e gravação são atômicas. & Concorrência não pode exceder a capacidade.\\
\end{tabelasrs}
\fonte{Elaboração didática deste guia.}
\end{quadrosrs}

A consulta pública independe de sessão. Criar, consultar registros pessoais e
cancelar exigem sessão válida. O direito de consultar e cancelar é verificado
pela titularidade, não apenas pelo conhecimento do identificador.
